\documentclass{article}
\usepackage{../math_template}

\author{Jonathan Kasongo}
\title{Simon Marais 2017 --- A1}

\begin{document}
\maketitle

\begin{problem}{Simon Marais 2017 --- A1}
The five sides and five diagonals of a regular pentagon are drawn on a piece of
paper. Two people play a game, in which they take turns to colour one of these ten
line segments. The first player colours line segments blue, while the second player
colours line segments red. A player cannot colour a line segment that has already
been coloured. A player wins if they are the first to create a triangle in their own
colour, whose three vertices are also vertices of the regular pentagon. The game is
declared a draw if all ten line segments have been coloured without a player winning.
\vspace{0.2cm}

Determine whether the first player, the second player, or neither player
can force a win.
\end{problem}

\begin{solution}{Solution}
The first player is indeed able to guarentee a win. \\

Step 1. The first player is able to colour two touching segments that lie
on the perimeter of the pentagon blue. Let us prove this statement.
On the first move player 1 should paint any of the 5 perimter segments
blue. Whatever player 2 does next, the segment that is painted blue has
two neighbour segments that touch it and also lie on the perimeter, so
at least one of them is left untouched after player 2's first move. Then
player 1 can safely colour one of the aformentioned neighbour segments.
\\

Step 2. Player 1 is able to force a win. Let us prove
this statement. Suppose player 2 has made their second move.
Consider the two diagonals $D_1$ and $D_2$ that touch the
vertex joining the two blue perimeter segments. \\

\textit{Case 1:} If one of $D_1, D_2$ is untouched, colour this diagonal
red. There are two mono-chromatic triangles that player 1 can make upon
their next turn, so no matter where player 2 goes they will lose.

\begin{center}
\begin{asy}
/* Geogebra to Asymptote conversion, documentation at artofproblemsolving.com/Wiki go to User:Azjps/geogebra */
import graph; size(160);
real labelscalefactor = 0.5; /* changes label-to-point distance */
pen dps = linewidth(0.7) + fontsize(10); defaultpen(dps); /* default pen style */
pen dotstyle = black; /* point style */
real xmin = -0.8675827508198834, xmax = 10.169539784533136, ymin = 0.3158377906479499, ymax = 9.437377114012188;  /* image dimensions */
pen qqccqq = rgb(0.,0.8,0.); pen ttttff = rgb(0.2,0.2,1.);
pair A = (2.,1.), B = (7.,1.), C = (8.545084971874736,5.755282581475766), D = (4.5,8.694208842938131);
/* draw figures */
draw(D--A );
draw(B--D,  qqccqq);
draw((0.45491502812526363,5.755282581475768)--C,  red); //
draw(A--C );
draw(B--(0.45491502812526363,5.755282581475768),  ttttff);
draw((0.45491502812526363,5.755282581475768)--D,  ttttff);
draw(D--C );
draw(C--B,  red); //
draw(B--A,  qqccqq);
draw(A--(0.45491502812526363,5.755282581475768),  ttttff);
/* dots and labels */
dot(A,dotstyle);
label("$A$", (2.040978521720044,1.1079987964585742), NE * labelscalefactor);
dot(B,dotstyle);
label("$B$", (7.0475859740501186,1.1079987964585742), NE * labelscalefactor);
dot(C,dotstyle);
label("$C$", (8.585490056940747,5.867643549346353), NE * labelscalefactor);
dot(D,dotstyle);
label("$D$", (4.544282247885081,8.808744788159084), NE * labelscalefactor);
dot((0.45491502812526363,5.755282581475768),dotstyle);
label("$E$", (0.5030744388294154,5.867643549346353), NE * labelscalefactor);
clip((xmin,ymin)--(xmin,ymax)--(xmax,ymax)--(xmax,ymin)--cycle);
/* end of picture */
\end{asy}
\end{center}

For example in this case, no matter what player 2 does next, player 1
is able to colour one of the green segments red and win the game. \\

\textit{Case 2:} If none of $D_1, D_2$ is untouched then player 1 can force
a win immediately. For example (referring to the above labelling)
if player 1 coloured $DE, AE$ blue and player 2 coloured $CE, BE$ red,
player 1 can colour $AD$ red and win. $\blacksquare$

\end{solution}

\end{document}
